\subsubsection*{EXP 2: Input-Dependent with $\tau$ annealing}

We investigated whether temperature annealing improves the InpDep model:

\begin{table}[h]
\centering
\caption{EXP~2: InpDep + $\tau$ annealing.}
\begin{tabular}{lcc}
\hline
\textbf{Configuration} & \textbf{Accuracy} & \textbf{Stability} \\
\hline
$\boldsymbol{\tau=0.5}$ \textbf{(fixed)} & $\boldsymbol{85.0\%}$ & \checkmark\ Stable \\
$\tau=0.1$ (fixed) & 83.7\% & \checkmark\ Stable \\
$\tau$: $1 \to 0.05$ (cosine) & 50.0\% (NaN) & $\times$ Explosion \\
$\tau$: $1 \to 0.05$ (linear) & 50.0\% (NaN) & $\times$ Explosion \\
\hline
\end{tabular}
\end{table}

\textbf{Aggressive annealing to $\tau=0.05$ causes numerical instability} in the Input-Dependent model. The hypernet generates log-potentials with high variance, and dividing by small $\tau$ causes overflow in $\exp()$. Fixed $\tau=0.5$ is optimal for InpDep. This problem motivates the transition to log-domain Sinkhorn (Section~\ref{sec:logdomain}).
